% =====================================================================
% EXAME DE QUALIFICAÇÃO DE DOUTORADO
% Yá-Sin Barcelos Mghazli — PEE/COPPE/UFRJ
%
% Base: template de apresentações do PEE (G. Gleizer), adaptado para
% 16:9 e organizado em módulos.
%
% Compilar:  make        (rascunho, com marcadores [TODO] visíveis)
%            make final  (versão de entrega, sem marcadores)
% =====================================================================

\documentclass[aspectratio=169,11pt]{beamer}

% --- Codificação e idioma -------------------------------------------
\usepackage[utf8]{inputenc}
\usepackage[T1]{fontenc}
\usepackage[brazil]{babel}
\usepackage{lmodern}

% --- Matemática e gráficos ------------------------------------------
\usepackage{amsmath,amssymb}
\usepackage{graphicx}
\usepackage{tikz}
\usepackage{booktabs}
\usepackage{multirow}
\usepackage{xcolor}

% Slides prontos vindos de fora (PowerPoint/Inkscape exportados em PDF)
% entram por aqui:  \includepdf[pages=2-4]{externos/diagrama.pdf}
\usepackage{pdfpages}

% Figuras: as da própria apresentação e, diretamente, as da tese.
\graphicspath{%
  {figuras/}%
  {imgs/}%
  {../Tese-de-Doutorado/figures/}%
  {../Tese-de-Doutorado/figures/experimentos/}%
  {../Tese-de-Doutorado/figures/tsad/}%
  {../Tese-de-Doutorado/figures/atlas/}%
}

% --- Identidade visual, macros e metadados --------------------------
% =====================================================================
% Identidade visual PEE/COPPE/UFRJ
% Derivado do template de G. Gleizer (tema Dresden + logos institucionais),
% adaptado para 16:9 e com posicionamento dos logos via TikZ.
% =====================================================================

\mode<presentation>{
  \usetheme{Dresden}
  \usecolortheme{seagull}
}

% O cabecalho de navegacao do Dresden e removido; a orientacao fica por
% conta do slide de agenda e dos separadores de secao.
\setbeamertemplate{headline}{}
\setbeamertemplate{navigation symbols}{}
\setbeamertemplate{frametitle}[default][center]

% ---------------------------------------------------------------------
% Paleta
% ---------------------------------------------------------------------
\definecolor{peeblue}{RGB}{1,123,165}
\definecolor{peedark}{RGB}{0,74,99}
\definecolor{peegray}{RGB}{95,105,110}
\definecolor{peeaccent}{RGB}{200,90,40}

\setbeamercolor*{palette primary}{fg=black,bg=peeblue!80}
\setbeamercolor*{palette secondary}{fg=black,bg=peeblue!80!gray!80}
\setbeamercolor*{palette tertiary}{fg=black,bg=peeblue!100}
\setbeamercolor*{palette quaternary}{fg=black,bg=peeblue!110}

\setbeamercolor{frametitle}{fg=peedark}
\setbeamercolor{title}{fg=peedark}
\setbeamercolor{structure}{fg=peeblue}
\setbeamercolor{alerted text}{fg=peeaccent}

% ---------------------------------------------------------------------
% Rodape
%
% O rodape de duas linhas do Dresden monta caixas cuja soma excede
% \paperwidth em 16:9, produzindo um Overfull \hbox por slide. Aqui ele e
% substituido por uma barra unica cujas larguras somam exatamente 1.
% ---------------------------------------------------------------------
\setbeamertemplate{footline}{%
  \leavevmode%
  \hbox{%
    \begin{beamercolorbox}[wd=.34\paperwidth,ht=2.4ex,dp=1.2ex,left]{author in head/foot}%
      \hspace*{1.2ex}\usebeamerfont{author in head/foot}\insertshortauthor
    \end{beamercolorbox}%
    \begin{beamercolorbox}[wd=.46\paperwidth,ht=2.4ex,dp=1.2ex,center]{title in head/foot}%
      \usebeamerfont{title in head/foot}\insertshorttitle
    \end{beamercolorbox}%
    \begin{beamercolorbox}[wd=.20\paperwidth,ht=2.4ex,dp=1.2ex,right]{date in head/foot}%
      \usebeamerfont{date in head/foot}\insertshortinstitute
      \hspace*{1.2em}\insertframenumber{}\hspace*{1.2ex}
    \end{beamercolorbox}}%
  \vskip0pt%
}

% ---------------------------------------------------------------------
% Tipografia e elementos de lista
% ---------------------------------------------------------------------
\setbeamerfont{frametitle}{size=\large,series=\bfseries}
\setbeamerfont{title}{size=\Large,series=\bfseries}
\setbeamertemplate{itemize item}{\color{peeblue}$\blacktriangleright$}
\setbeamertemplate{itemize subitem}{\color{peegray}\textendash}
\setbeamertemplate{enumerate item}{\color{peeblue}\insertenumlabel.}
\setbeamercolor{caption name}{fg=peedark}
\setbeamerfont{caption}{size=\scriptsize}

% Blocos com cabecalho institucional
\setbeamertemplate{blocks}[rounded][shadow=false]
\setbeamercolor{block title}{fg=white,bg=peeblue}
\setbeamercolor{block body}{bg=peeblue!8}
\setbeamercolor{block title alerted}{fg=white,bg=peeaccent}
\setbeamercolor{block body alerted}{bg=peeaccent!8}

% ---------------------------------------------------------------------
% Logos de fundo (TikZ: posicionamento absoluto, imune ao aspect ratio)
% ---------------------------------------------------------------------
\newcommand{\peelogomarca}{%
  \begin{tikzpicture}[remember picture,overlay]
    \node[anchor=north west,inner sep=0pt]
      at ([shift={(0.28cm,-0.22cm)}]current page.north west)
      {\includegraphics[height=0.085\paperheight,keepaspectratio]{imgs/pee-logo-short.png}};
  \end{tikzpicture}%
}

\newcommand{\peelogoscapa}{%
  \begin{tikzpicture}[remember picture,overlay]
    \node[anchor=north west,inner sep=0pt]
      at ([shift={(0.9cm,-0.7cm)}]current page.north west)
      {\includegraphics[height=0.135\paperheight,keepaspectratio]{imgs/pee-logo.png}};
    \node[anchor=north east,inner sep=0pt]
      at ([shift={(-0.9cm,-0.7cm)}]current page.north east)
      {\includegraphics[height=0.135\paperheight,keepaspectratio]{imgs/coppe-logo.pdf}};
  \end{tikzpicture}%
}

% A marca dagua discreta acompanha todos os slides de conteudo.
\usebackgroundtemplate{\peelogomarca}

% A area util recua nas duas margens: o titulo centralizado fica livre da
% marca dagua no canto superior esquerdo sem estourar a caixa.
\setbeamersize{text margin left=1.1em,text margin right=1.1em}

% ---------------------------------------------------------------------
% Separador de secao automatico
% ---------------------------------------------------------------------
\AtBeginSection[]{%
  \begin{frame}[plain,noframenumbering]
    \vfill
    \centering
    \begin{beamercolorbox}[sep=10pt,center,shadow=false,rounded=true]{title}
      \usebeamerfont{title}\color{peedark}\insertsectionhead\par
    \end{beamercolorbox}
    \vfill
  \end{frame}%
}

% =====================================================================
% Macros de apoio aos slides
% =====================================================================

% Marcador visível de conteúdo pendente. Some ao compilar em modo final
% ("make final"), para não ir pendência nenhuma na versão de entrega.
\newif\ifrascunho
\rascunhotrue
\ifdefined\modofinal\rascunhofalse\fi

\newcommand{\pendente}[1]{%
  \ifrascunho\textcolor{peeaccent}{\footnotesize\textbf{[TODO: #1]}}\fi}

% Destaque semântico.
\newcommand{\dest}[1]{\textcolor{peedark}{\textbf{#1}}}
\newcommand{\alerta}[1]{\textcolor{peeaccent}{\textbf{#1}}}

% Figura ocupando o slide inteiro, com legenda opcional no rodapé.
% Uso: \framefigura{caminho/figura}{Legenda}
\newcommand{\framefigura}[2]{%
  \begin{frame}[plain]
    \centering
    \includegraphics[width=0.92\paperwidth,height=0.82\paperheight,keepaspectratio]{#1}\\[0.3em]
    {\scriptsize\color{peegray}#2}
  \end{frame}}

% Slide de dois blocos lado a lado.
% Uso: \duascolunas{esquerda}{direita}
\newcommand{\duascolunas}[2]{%
  \begin{columns}[T,onlytextwidth]
    \begin{column}{0.48\textwidth}#1\end{column}
    \begin{column}{0.48\textwidth}#2\end{column}
  \end{columns}}

% Abreviações recorrentes do domínio.
\newcommand{\TSAD}{\mbox{TSAD}}
\newcommand{\DQM}{\mbox{DQM}}
\newcommand{\LAr}{\mbox{LAr}}

% =====================================================================
% Metadados do Exame de Qualificação
% Extraídos de ../Tese-de-Doutorado/main.tex — manter sincronizado.
% =====================================================================

% Título curto para o rodapé; título completo na capa.
\title[Sistema Multi-Agente para DQM no ATLAS]{%
  Um Sistema Multi-Agente Colaborativo para o Monitoramento da
  Qualidade de Dados do Calorímetro LAr do ATLAS via Detecção de
  Anomalias, Reconstrução e Previsão de Séries Temporais}

\subtitle{Exame de Qualificação de Doutorado}

\author[Yá-Sin B. Mghazli]{Yá-Sin Barcelos Mghazli}

\institute[UFRJ/COPPE/PEE]{%
  Programa de Engenharia Elétrica\\
  COPPE --- Universidade Federal do Rio de Janeiro}

% TODO: confirmar a data exata com a secretaria do PEE.
\date{Rio de Janeiro, setembro de 2026}

% ---------------------------------------------------------------------
% Orientação e banca (usados no slide de abertura)
% ---------------------------------------------------------------------
\newcommand{\orientadores}{%
  Prof. José Manoel de Seixas, D.Sc.\\
  Prof. Juan Lieber Marin, D.Sc.}

\newcommand{\bancaexaminadora}{%
  Prof. José Manoel de Seixas, D.Sc.\\
  Prof. Juan Lieber Marin, D.Sc.\\
  Prof. José Gabriel Rodríguez Carneiro Gomes, D.Sc.\\
  Prof. Edmar Egidio Purcino de Souza, D.Sc.\\
  Prof. Tiago Alessandro Espinola Ferreira, D.Sc.}


% =====================================================================
\begin{document}

% =====================================================================
% Capa, identificação da banca e agenda
% =====================================================================

% A capa usa os dois logos ampliados e suprime a marca d'água discreta.
{
\usebackgroundtemplate{\peelogoscapa}
\begin{frame}[plain,noframenumbering]
  % O recuo vertical e a medida reduzida mantêm o título abaixo da faixa
  % ocupada pelos logos, que sangram até ~0.20\paperheight.
  \vspace{2.6em}
  \begin{center}
    \begin{minipage}{0.86\textwidth}
      \centering
      {\large\bfseries\color{peedark}\inserttitle\par}
    \end{minipage}\par
    \vspace{0.9em}
    {\color{peegray}\normalsize\insertsubtitle\par}
    \vspace{1.4em}
    {\large\insertauthor\par}
    \vspace{0.4em}
    {\footnotesize\color{peegray}\insertinstitute\par}
    \vspace{1.0em}
    {\footnotesize\insertdate\par}
  \end{center}
\end{frame}
}

\begin{frame}[noframenumbering]{Orientação e banca examinadora}
  \begin{columns}[T,onlytextwidth]
    \begin{column}{0.42\textwidth}
      \textbf{\color{peedark}Orientadores}
      \vspace{0.5em}

      {\footnotesize\orientadores}
    \end{column}
    \begin{column}{0.54\textwidth}
      \textbf{\color{peedark}Banca examinadora}
      \vspace{0.5em}

      {\footnotesize\bancaexaminadora}
    \end{column}
  \end{columns}
  \vfill
  {\scriptsize\color{peegray}
   Programa de Engenharia Elétrica --- COPPE/UFRJ}
\end{frame}

\begin{frame}{Agenda}
  \tableofcontents[hideallsubsections]
\end{frame}


% =====================================================================
% PARTE I — Contexto, problema e objetivos      [~8 min]
% Fonte: Chapters/introducao, Chapters/contexto_cern_atlas
% =====================================================================

\section{Contexto e Problema}

\begin{frame}{O experimento ATLAS e o calorímetro LAr}
  \duascolunas{%
    \begin{itemize}
      \item Colisões a \dest{40\,MHz}, vazão da ordem de \dest{70\,TB/s}.
      \item O calorímetro de argônio líquido (\LAr) fornece as grandezas
            usadas na reconstrução de energia.
      \item O monitoramento de qualidade de dados (\DQM) decide o que é
            aproveitável para física.
    \end{itemize}
    \pendente{escolher a figura do detector --- ver figures/atlas/}
  }{%
    \centering
    \pendente{\texttt{\textbackslash includegraphics} da figura do LAr}
  }
\end{frame}

\begin{frame}{O gargalo do monitoramento de qualidade de dados}
  \begin{itemize}
    \item Inspeção majoritariamente \dest{manual} e por limiares fixos.
    \item O critério Toast e suas limitações.
      \pendente{enunciar as três limitações do Cap. 3}
    \item Defeitos D3--D11: heterogêneos em escala, duração e assinatura.
    \item Consequência: latência de diagnóstico e perda de dados válidos.
  \end{itemize}
\end{frame}

\begin{frame}{Questão de pesquisa e hipótese}
  \begin{block}{Questão}
    \pendente{transcrever a questão de pesquisa do Cap. 1}
  \end{block}
  \vspace{0.5em}
  \begin{alertblock}{Hipótese}
    \pendente{transcrever a hipótese central --- a sinergia entre detecção,
    reconstrução e previsão supera os agentes isolados}
  \end{alertblock}
\end{frame}

\begin{frame}{Objetivos}
  \textbf{\color{peedark}Objetivo geral}
  \begin{itemize}
    \item \pendente{objetivo geral do Cap. 1}
  \end{itemize}
  \vspace{0.6em}
  \textbf{\color{peedark}Objetivos específicos}
  \begin{enumerate}
    \item \pendente{OE1}
    \item \pendente{OE2}
    \item \pendente{OE3}
  \end{enumerate}
\end{frame}

% =====================================================================
% PARTE I — Fundamentação                       [~5 min]
% Fonte: Chapters/fundamentacao
% =====================================================================

\section{Fundamentação}

\begin{frame}{Séries temporais: representação e pressupostos}
  \begin{itemize}
    \item Três casos de observação: amostragem uniforme, sequência ordenada
          por índice e série de eventos.
    \item O que a tese \alerta{não} assume.
      \pendente{listar as delimitações do Cap. 2}
    \item Estacionariedade, regime e deriva: distinção operacional em
          relação a anomalia.
  \end{itemize}
\end{frame}

\begin{frame}{Sinal, ruído, resíduo e anomalia}
  \begin{block}{Distinção central}
    Resíduo elevado \alerta{não implica} anomalia.
  \end{block}
  \vspace{0.4em}
  \begin{itemize}
    \item Decomposição aditiva e multiplicativa.
      \pendente{trazer eq:aditivo e eq:multiplicativo do Cap. 2}
    \item Estatísticas estimadas somente em partições autorizadas.
  \end{itemize}
\end{frame}

\begin{frame}{Qualidade de dados: ausência, corrupção e integridade}
  \begin{itemize}
    \item Mecanismos de ausência: \dest{MCAR}, \dest{MAR}, \dest{MNAR}.
    \item Implicação sobre a validade da imputação.
    \item \pendente{ligar aos requisitos derivados no Apêndice B}
  \end{itemize}
\end{frame}

% =====================================================================
% PARTE II — Revisão crítica do estado da arte  [~10 min]
% Fonte: Chapters/tsad, Chapters/reconstruction, Chapters/forecasting
% =====================================================================

\section{Estado da Arte}

\subsection{Detecção de anomalias}

\begin{frame}{Taxonomia de anomalias em séries temporais}
  \duascolunas{%
    \begin{itemize}
      \item Pontuais
      \item Contextuais
      \item Coletivas
    \end{itemize}
    \vspace{0.5em}
    {\scriptsize\color{peegray}Critérios em prosa; sem equações ilustrativas.}
  }{%
    \centering
    \pendente{usar \texttt{anomaly\_taxonomy.png} do workspace}
  }
\end{frame}

\begin{frame}{Famílias de métodos e o problema da avaliação}
  \begin{itemize}
    \item Distância, densidade e predição; realizações neurais.
    \item \dest{Benchmarks} e o viés de \emph{point-adjust}.
    \item Medidas de localização: Tatbul, VUS, afiliação, PATE.
    \item \pendente{fechar com a lacuna que justifica a proposta}
  \end{itemize}
\end{frame}

\subsection{Reconstrução e previsão}

\begin{frame}{Reconstrução: contrato da tarefa}
  \begin{itemize}
    \item Três estatutos do alvo.
      \pendente{Cap. 5}
    \item Reconstrução \alerta{$\neq$} correção de anomalia.
  \end{itemize}
\end{frame}

\begin{frame}{Previsão: do ARIMA aos Transformers}
  \begin{itemize}
    \item Métodos clássicos, aprendizado de máquina e profundo.
    \item Transformers para séries temporais: ganho real e custo.
    \item \pendente{Cap. 6}
  \end{itemize}
\end{frame}

\begin{frame}{Síntese: as três lacunas}
  \begin{enumerate}
    \item \pendente{lacuna 1}
    \item \pendente{lacuna 2}
    \item \pendente{lacuna 3}
  \end{enumerate}
  \vfill
  \begin{block}{Consequência}
    \pendente{a transição para a proposta}
  \end{block}
\end{frame}

% =====================================================================
% PARTE III — Decisão metodológica e contribuição  [~12 min]
% Fonte: Chapters/estrategias_fusao, Chapters/sistema_multi_agente
% Núcleo da apresentação: é aqui que a banca decide se a tese se sustenta.
% =====================================================================

\section{Proposta}

\begin{frame}{Da matriz de decisão à arquitetura}
  \begin{itemize}
    \item Critérios de seleção metodológica.
      \pendente{matriz de decisão do Cap. 7}
    \item Por que multi-agente, e não um modelo monolítico.
  \end{itemize}
\end{frame}

% A arquitetura merece o slide inteiro: é a figura que a banca vai olhar
% mais tempo. O TikZ da tese já existe e exporta em vetorial.
\begin{frame}[plain]
  \centering
  \pendente{arquitetura\_multiagente.tex --- compilar standalone e incluir o PDF}
  % \includegraphics[width=0.92\paperwidth,height=0.84\paperheight,keepaspectratio]%
  %   {figuras/arquitetura_multiagente.pdf}
\end{frame}

\begin{frame}{Estratégias de fusão}
  \duascolunas{%
    \textbf{\color{peedark}Nível de score}
    \begin{itemize}\footnotesize
      \item \pendente{}
    \end{itemize}
  }{%
    \textbf{\color{peedark}Nível de decisão}
    \begin{itemize}\footnotesize
      \item \pendente{}
    \end{itemize}
  }
\end{frame}

\begin{frame}{Coordenação: intra-grupo, hierárquica e blackboard}
  \begin{itemize}
    \item Coordenação intra-grupo entre agentes da mesma tarefa.
    \item Coordenação hierárquica entre os três grupos.
    \item \dest{Blackboard} e meta-coordenação.
    \item \pendente{Cap. 8}
  \end{itemize}
\end{frame}

\begin{frame}{Hipóteses experimentais}
  \begin{enumerate}
    \item \pendente{H1}
    \item \pendente{H2}
    \item \pendente{H3 --- a hipótese de sinergia, verificada por ablação}
  \end{enumerate}
\end{frame}

% =====================================================================
% PARTE IV — Validação experimental             [~8 min]
% Fonte: Chapters/experimentos
% =====================================================================

\section{Resultados Preliminares}

\begin{frame}{Protocolo experimental}
  \begin{itemize}
    \item \textbf{Dados:} \pendente{datasets e o gerador do Apêndice D}
    \item \textbf{Baselines:} \pendente{}
    \item \textbf{Métricas:} \pendente{justificar a escolha, dado o Cap. 4}
  \end{itemize}
\end{frame}

\begin{frame}{Agentes individuais}
  \centering
  \pendente{tabela ou figura por família --- ver figures/experimentos/}
\end{frame}

\begin{frame}{Sistema integrado e ablação}
  \begin{block}{Prova de sinergia}
    \pendente{o resultado que sustenta H3}
  \end{block}
  \vspace{0.4em}
  \pendente{gráfico de ablação}
\end{frame}

\begin{frame}{Aplicação em dados reais do ATLAS}
  \centering
  \pendente{estudo de caso em dados reais --- figures/atlas/}
\end{frame}

% =====================================================================
% Cronograma, contribuições e encerramento      [~5 min]
% Fonte: Chapters/conclusoes
% =====================================================================

\section{Cronograma e Trabalhos Futuros}

\begin{frame}{O que já está feito}
  \begin{itemize}
    \item \pendente{revisão crítica das três tarefas --- Caps. 4--6}
    \item \pendente{formulação da arquitetura --- Caps. 7--8}
    \item \pendente{validação parcial --- Cap. 9}
  \end{itemize}
\end{frame}

\begin{frame}{O que falta e em que prazo}
  \centering
  \pendente{tabela de cronograma por semestre até a defesa}
  % Sugestão: tabular com colunas Atividade | 2026.2 | 2027.1 | 2027.2
\end{frame}

\begin{frame}{Contribuições esperadas}
  \begin{enumerate}
    \item \pendente{C1}
    \item \pendente{C2}
    \item \pendente{C3}
  \end{enumerate}
\end{frame}

\begin{frame}[plain]
  \vfill
  \centering
  {\Large\color{peedark}\textbf{Obrigado.}}\\[1.2em]
  {\footnotesize\color{peegray}
   Yá-Sin Barcelos Mghazli\\
   Programa de Engenharia Elétrica --- COPPE/UFRJ}
  \vfill
\end{frame}


% =====================================================================
% Slides de apoio (após o "Obrigado", fora da numeração principal)
%
% Regra prática: um slide de backup por pergunta que você consegue
% antecipar. É aqui que a qualificação é ganha ou perdida.
% =====================================================================

\appendix

\begin{frame}[noframenumbering]{Apoio --- índice}
  \begin{itemize}\footnotesize
    \item \pendente{formulações matemáticas --- Apêndice C}
    \item \pendente{hiperparâmetros e proveniência das execuções}
    \item \pendente{requisitos derivados por domínio --- Apêndice B}
    \item \pendente{inventário taxonômico de modelos --- Apêndice A}
    \item \pendente{resultados complementares D7/D8 --- Apêndice D}
  \end{itemize}
\end{frame}


\end{document}
